% GENERATED by qa/build_cjk_reader_streams.py; DO NOT EDIT.
% Continuous preliminary matter with complete source milestones.
\VerdierSemanticAnchor{front-title}
\begin{titlepage}
  \VerdierSourcePageBlock{0002}
  \centering
  {\Huge\bfseries 239\hfill ASTÉRISQUE\par}
  \vspace{1.5em}
  {\LARGE\bfseries 1996\par}
  \vfill
  {\Huge\bfseries アーベル圏の\par}
  \vspace{0.4em}
  {\Huge\bfseries 導来圏\par}
  \vspace{3em}
  {\Large\bfseries Jean-Louis VERDIER\par}
  \vfill
  {\Large\bfseries フランス数学会\par}
  \smallskip
  {\small\bfseries フランス国立科学研究センターの助成により刊行\par}
\end{titlepage}

\clearpage
\VerdierSemanticAnchor{front-ams-classification}
\VerdierSourcePageBlock{0003}
\VerdierReaderChapterStar{AMS 主題分類}
\addcontentsline{toc}{chapter}{AMS 主題分類}

14F20, 18E30, 18E35, 18G05, 18G10, 18G15, 18G20, 18G35,
18G40, 55U15, 55U25, 55U30.

\clearpage
\VerdierSemanticAnchor{front-portrait}
\VerdierSourcePageBlock{0004}
\thispagestyle{empty}
\vspace*{0.62\textheight}
\begin{center}
Jean-Louis VERDIER、1977年
\end{center}

\VerdierSourcePageSilent{0005}%
\clearpage
\VerdierSemanticAnchor{front-preface}
\VerdierSourcePageBlock{0006}
\VerdierReaderChapterStar{序文}
\addcontentsline{toc}{chapter}{序文}

1989年8月25日に悲劇的な死を遂げた Jean-Louis Verdier は、1967年6月14日、
パリにおいて『アーベル圏の導来圏』と題する国家博士論文の審査を受けた。
本書のテキストは、この論文のタイプ原稿の最終版である。編者が行ったのは細部
の誤りを訂正することに限られる。学部に提出されていたのは「序論」と題する
部分だけであり、本文は審査時にも未完で、その後も未完のままであった。

この仕事の成立事情はかなりよく知られている（\emph{参照} [I]）。
Grothendieck は1960年代初頭、1958年のエディンバラ国際数学者会議での講演を
受けて構想していた双対性定理の大規模な一般化に必要なホモロジー代数の枠組み
を与えるため、導来圏という着想を得た。彼は Verdier に、その形式体系の構築
を博士論文の課題として提案した。1963年、Verdier は理論の本質的内容を要約
した成果集『導来圏、いくつかの結果（状態0）』を IHÉS の謄写版として刊行した
（[SGA 4 1/2] に再録）。導来圏の概念は Grothendieck に由来するが、それに
先立って三角圏の概念を導入したのは Verdier である。安定ホモトピー論に着想
を得たそこでの公理的枠組み（\emph{参照} [P]）は、驚くほど豊かな成果をもたらす
ことになった。

Verdier が自らの仕事を包括的に記述し始めたのは、それから何年も後のことで
ある。なぜ完成させなかったのか。多くの仮説が唱えられてきたが、私は次の説明
に傾く。欠けている部分は導来関手の理論である。Verdier がそこまで書き進めた
頃には、導来圏はすでに代数幾何で広く使われ、さまざまな変形と一般化が発展
していた。Deligne 流の各点での導来可能性、非加法的導来関手（Quillen 流の
ホモトピー代数）、フィルター付き導来関手などである。その時点で、成果集が
定めた枠組みは不十分になっていた。適切な観点を見いだすには新たな基礎研究が
必要だったが、当時ほかの問題に関心を向けていた Verdier には、その余裕も、
おそらく意欲もなかったのだと私は思う。時がたつにつれ導来圏に関する論文が
現れ、その利用は広がり、形式体系は数学の別分野（たとえば線形偏微分方程式）
にも浸透した。こうして Verdier の目には、この仕事を完成させることがますます
不要に映るようになったのである……。

\VerdierSourcePageBlock{0007}
未完ではあるものの、Verdier のこのテキストは貴重で、まさに代えがたい参考
文献である（たとえば、ここに示されたスペクトル系列の独創的で優美な提示は、
私の知る限りほかのどこにも見られない）。それがついに日の目を見たことを、
心から喜びたい。

\begin{flushright}
パリ、1996年4月27日\\
Luc Illusie
\end{flushright}

\section*{参考文献}

\begin{description}[labelwidth=6.2em,leftmargin=6.8em]
\item[{[I]}] L. Illusie, \emph{Catégories dérivées et dualité, travaux de
J.-L. Verdier}, L'Enseignement Mathématique, 36 (1990), 369--391.
（1989年10月19日、Université Paris VII 主催の Jean-Louis Verdier 追悼式
における講演。）

\item[{[P]}] D. Puppe, \emph{On the formal structure of stable homotopy
theory}, Colloquium on Algebraic Topology, Aarhus Universitet (1962),
65--71.

\item[{\mbox{[SGA 4 1/2]}}] ``Cohomologie étale'', Séminaire de Géométrie
Algébrique du Bois-Marie, SGA 4 1/2, P. Deligne, Lecture Notes in
Mathematics 569, Springer-Verlag (1977), 262--311.
\end{description}
\clearpage
\VerdierSemanticAnchor{front-editors-note}
\VerdierSourcePageBlock{0008}
\VerdierReaderChapterStar{編者注}
\addcontentsline{toc}{chapter}{編者注}

\VerdierSemanticAnchor{front-editors-note-publication-history}
J.-L. Verdier を追悼して Université Paris VII が1989年10月19日に催した式典
において、M. Artin、P. Cartier、L. Illusie は J.-L. Verdier の博士論文を
出版する着想を得た。2年後、私も彼らに加わり、J. Le Potier と M. Broué の
支援を得て、本文を \TeX{} に入力するための資金を Paris VII（URA 212）と
ÉNS（DMI）から確保した。入力完了後、その結果を Verdier のタイプ原稿と入念
に照合した。しかし Verdier 自身が最終訂正を行っていなかったため、数学的な
校閲がなお必要であった。

\VerdierSemanticAnchor{front-editors-note-chapter-i-modifications}
この精密な校閲は、P. Cartier（第I章）、B. Keller（第III章）、そして全文に
ついては M. Zisman と私が、細心の注意を払って行った。列挙するほどでもない
多数の誤植を除けば、主な変更は次のとおりである。第I章の命題 (3.1.5) は同値
を主張していた。この逆向き（後では一度も使われない）が偽なので、同値を含意
に置き換えた。命題 (3.4.1) には符号の誤りがあり、第III章 ((1.3.4)、(b)、
(3.2.7.3)) の符号の誤りを引き起こしていた。そのため「反対三角圏」
（第II章、1.1.7）の定義を調整し、次の主張が真であり続けるようにした：
\emph{反対圏の導来圏は、元の導来圏の反対三角圏である}。

\VerdierSemanticAnchor{front-editors-note-chapter-ii-modifications}
第II章では、§1.2 の命題の順序を変更した。Verdier は執筆後、二つの完全三角
の直和が完全三角になるという主張は三角圏の他の公理から従うことに気づいて
いた。そのため公理の一覧からは削除したが、それに伴う再配列を行っていなかった
のである。さらに、後で使われず、検証にも失敗した段落 (2.3.6) と (2.3.7) を
削除した。定常スペクトル対象 (4.4.2) の定義はわずかに修正した。元の定義では
弱すぎて、後続の定理 (4.4.3) を証明できなかったためである。この変更が以後
の内容に影響しないことも確認した。

\VerdierSourcePageBlock{0009}
\VerdierSemanticAnchor{front-editors-note-apparatus}
記号索引と用語索引は補完した。一方、参考文献は変更しなかった（ただし SGA
各巻については ``Lecture Notes'' での確定書誌情報を示した）。とりわけ、
参考文献番号13の欠落は補っていない。

\VerdierSemanticAnchor{front-editors-note-acknowledgements}
この企画の実現をさまざまな形で助けてくださった M. Broué、A. Bruguières、
P. Cartier、L. Gruson、L. Illusie、B. Keller、J. Le Potier、M. Zisman、
ならびに Université Paris VII、École Normale Supérieure、Institut des
Hautes Études Scientifiques に心から感謝したい。また、本書を辛抱強く
\textsc{AMS}-\TeX{} に入力してくださった C. Defosse と C. Gourgues にも
感謝する。

\VerdierSemanticAnchor{front-editors-note-memorial}
1994年12月に早逝した L. Doustaing にも、この感謝を伝えたかった。

\VerdierSemanticAnchor{front-editors-note-signature}
\begin{flushright}
パリ、1996年6月14日\\
Georges Maltsiniotis
\end{flushright}
\clearpage
\VerdierSemanticAnchor{front-publications}
\VerdierSourcePageBlock{0010}
\VerdierReaderChapterStar{J.-L. Verdier 著作目録}
\addcontentsline{toc}{chapter}{J.-L. Verdier 著作目録}

\small
\begin{description}
\VerdierSemanticAnchor{front-publications-01}
\item[{[01]}] \emph{保型形式に付随する積分について}（G. Shimura による）、
Séminaire Bourbaki 1960--61, no. 216, pp. 216-1--216-27,
Benjamin (1966).

\VerdierSemanticAnchor{front-publications-02}
\item[{[02]}] \emph{Poincaré 双対定理}, C.R.A.S. vol. 256,
pp. 2084--2086, Gauthier-Villars (1963).

\VerdierSemanticAnchor{front-publications-03}
\item[{[03]}] \emph{数体のエタール・コホモロジーに関するセミナー}
（M. Artin との共著）, Woods Hole Conference on Algebraic Geometry,
A.M.S. (1964).

\VerdierSemanticAnchor{front-publications-04}
\item[{[04]}] \emph{射影系の本質的同値}, C.R.A.S. vol. 261,
pp. 4950--4953, Gauthier-Villars (1965).

\VerdierSemanticAnchor{front-publications-05}
\item[{[05]}] \emph{局所コンパクト空間の次元}, C.R.A.S. vol. 261,
pp. 5293--5296, Gauthier-Villars (1965).

\VerdierSemanticAnchor{front-publications-06}
\item[{[06]}] \emph{副有限群のコホモロジーにおける双対性}, Jean-Pierre
Serre, ``Galois Cohomology'' の付録, Lecture Notes in Mathematics 5,
pp. 183--206, Springer (1964), 第4版 (1973).

\VerdierSemanticAnchor{front-publications-07}
\item[{[07]}] \emph{局所コンパクト空間上の構成可能層}, C.R.A.S. vol. 262,
pp. 12--15, Gauthier-Villars (1966).

\VerdierSemanticAnchor{front-publications-08}
\item[{[08]}] \emph{局所コンパクト空間のコホモロジーにおける双対性},
Séminaire Bourbaki 1965--66, no. 300, pp. 300-01--300-13,
Benjamin (1966).

\VerdierSemanticAnchor{front-publications-09}
\item[{[09]}] \emph{エタール・コホモロジーにおける Lefschetz 不動点公式},
Conference on Local Fields, 1966年 Driebergen の NUFFIC Summer School,
pp. 199--214, Springer (1967).

\VerdierSemanticAnchor{front-publications-10}
\item[{[10]}] \emph{スキームのエタール・コホモロジーにおける双対定理},
Conference on Local Fields, 1966年 Driebergen の NUFFIC Summer School,
pp. 184--198, Springer (1967).

\VerdierSemanticAnchor{front-publications-11}
\item[{[11]}] \emph{連接層のねじれ逆像に対する基底変換}, Bombay
Colloquium on Algebraic Geometry in 1968, pp. 393--408, Oxford
University Press (1969).
\end{description}
\normalsize
\VerdierSourcePageContinue{0011}
\small
\begin{description}
\VerdierSemanticAnchor{front-publications-12}
\item[{[12]}] \emph{局所コンパクト空間のコホモロジーに対する双対定理},
Poincaré Duality, Heidelberg-Strasbourg Seminar 1966--67, exp. 4,
Publ. I.R.M.A. Strasbourg, no. 3, pp. 1--36 (1969).

\VerdierSemanticAnchor{front-publications-13}
\item[{[13]}] \emph{核定理}, Poincaré Duality, Heidelberg-Strasbourg
Seminar 1966--67, exp. 9, Publ. I.R.M.A. Strasbourg, no. 3,
pp. 1--52 (1969).

\VerdierSemanticAnchor{front-publications-14}
\item[{[14]}] \emph{Wilder の定理について}, Applications of categorical
algebra, Proc. Symp. Pure Math. 17, pp. 184--191, A.M.S. (1970).

\VerdierSemanticAnchor{front-publications-15}
\item[{[15]}] \emph{Grauert の連接性定理の簡明な証明}（R. Kiehl との共著）, 
Math. Ann. 195, pp. 24--50, Springer (1971).

\VerdierSemanticAnchor{front-publications-16}
\item[{[16]}] \emph{連接コホモロジーをもつ解析層複体のコホモロジー空間上の
位相}, Bull. Soc. Math. France 99, pp. 337--343, Gauthier-Villars
(1971).

\VerdierSemanticAnchor{front-publications-17}
\item[{[17]}] \emph{複素解析幾何における相対双対性}（J.-P. Ramis、
G. Ruget との共著）, Inv. Math. 13, pp. 261--283, Springer (1971).

\VerdierSemanticAnchor{front-publications-18}
\item[{[18]}] S.G.A. 4 (1963--64)「トポスの理論とスキームのエタール・
コホモロジー」：\emph{前層}（A. Grothendieck との共著）, pp. 1--184；
\emph{位相と層}, pp. 219--263；\emph{層の圏の関手性}, pp. 265--297；
\emph{トポス}（A. Grothendieck との共著）, pp. 299--519,
Lecture Notes in Mathematics 269, Springer (1972).

\VerdierSemanticAnchor{front-publications-19}
\item[{[19]}] S.G.A. 4 (1963--64)「トポスの理論とスキームのエタール・
コホモロジー」：\emph{トポスにおけるコホモロジー}, pp. 1--62；
\emph{有限性条件、ファイバー化トポスとサイト、および極限移行問題への応用}
（A. Grothendieck との共著）, pp. 163--336, Lecture Notes in
Mathematics 270, Springer (1972).

\VerdierSemanticAnchor{front-publications-20}
\item[{[20]}] \emph{Euler-Poincaré 標数}, Bull. Soc. Math. France 101,
pp. 441--445, Gauthier-Villars (1973).

\VerdierSemanticAnchor{front-publications-21}
\item[{[21]}] \emph{Le Potier の定理}, Seminar on Analytic Geometry
1972--73, Various aspects of positivity, Astérisque 17, pp. 68--78,
S.M.F. (1974).

\VerdierSemanticAnchor{front-publications-22}
\item[{[22]}] \emph{$l$ 進コホモロジーの Frobenius 自己準同型の特性多項式
の $l$ からの独立性}（P. Deligne による）, Séminaire Bourbaki
1972--73, no. 423, Lecture Notes in Mathematics 383, pp. 98--115,
Springer (1974).

\VerdierSemanticAnchor{front-publications-23}
\item[{[23]}] \emph{Riemann 幾何におけるいくつかの交叉問題}
（P. Marry との共著）, J. Diff. Geom. 12, pp. 345--376, A.M.S. (1977).
\end{description}
\normalsize
\VerdierSourcePageContinue{0012}
\small
\begin{description}
\VerdierSemanticAnchor{front-publications-24}
\item[{[24]}] \emph{Whitney 層化と Bertini-Sard 定理}, Inv. Math. 36,
pp. 295--312, Springer (1976).

\VerdierSemanticAnchor{front-publications-25}
\item[{[25]}] \emph{サイクルに付随するホモロジー類}, Seminar on
Analytic Geometry 1974--75, Astérisque 36--37, pp. 101--151,
S.M.F. (1976).

\VerdierSemanticAnchor{front-publications-26}
\item[{[26]}] \emph{完全交叉に対する Riemann-Roch 定理}, Seminar on
Analytic Geometry 1974--1975, Astérisque 36--37, pp. 189--228,
S.M.F. (1976).

\VerdierSemanticAnchor{front-publications-27}
\item[{[27]}] \emph{特異点をもちうる代数多様体に対する Riemann-Roch 定理}
（P. Baum、W. Fulton、R. MacPherson による）, Séminaire Bourbaki
1974--75, no. 464, Lecture Notes in Mathematics 514, pp. 159--175,
Springer (1976).

\VerdierSemanticAnchor{front-publications-28}
\item[{[28]}] \emph{導来圏}（1963年執筆、1976年刊行）, Seminar on
Algebraic Geometry 4 1/2, Étale Cohomology, Lecture Notes in Mathematics
569, pp. 262--311, Springer (1977).

\VerdierSemanticAnchor{front-publications-29}
\item[{[29]}] \emph{代数微分方程式}, Séminaire Bourbaki 1977--78,
no. 512, Lecture Notes in Mathematics 710, pp. 101--122, Springer (1979).

\VerdierSemanticAnchor{front-publications-30}
\item[{[30]}] \emph{インスタントン}, The Yang-Mills equations,
É.N.S. Seminar 1977--78, Astérisque 71--72, pp. 105--134, S.M.F. (1980).

\VerdierSemanticAnchor{front-publications-31}
\item[{[31]}] \emph{Chern 類の特殊化}, Euler-Poincaré characteristic,
É.N.S. Seminar 1978--79, Astérisque 82--83, pp. 149--159, S.M.F. (1981).

\VerdierSemanticAnchor{front-publications-32}
\item[{[32]}] \emph{Lie 代数、Hamilton 系、代数曲線}（M. Adler と
P. Van Moerbeke による）, Séminaire Bourbaki 1980--81, no. 566,
Lecture Notes in Mathematics 901, pp. 85--94, Springer (1981).

\VerdierSemanticAnchor{front-publications-33}
\item[{[33]}] \emph{アフィン Lie 代数の表現：物理学のいくつかの問題への応用}
（E. Date、M. Jimbo、M. Kashiwara、T. Miwa による）, Séminaire
Bourbaki 1981--82, no. 596, Astérisque 92--93, pp. 365--377,
S.M.F. (1982).

\VerdierSemanticAnchor{front-publications-34}
\item[{[34]}] \emph{行列 $S$ の構成}（A. Treibich-Kohn との共著）,
Mathematics and Physics, É.N.S. Seminar 1979--82, Progress in Mathematics
37, pp. 175--196, Birkhäuser (1983).

\VerdierSemanticAnchor{front-publications-35}
\item[{[35]}] \emph{2次元 $\sigma$ 模型と $S^2$ から $S^{2n}$ への調和写像},
Group Theoretical Methods in Physics, Lecture Notes in Physics 180,
pp. 136--141, Springer (1983).

\VerdierSemanticAnchor{front-publications-36}
\item[{[36]}] \emph{超局所幾何}, Algebraic Geometry, Japan-France
Conference, Tokyo-Kyoto 1982, Lecture Notes in Mathematics 1016,
pp. 125--133, Springer (1983).
\end{description}
\normalsize
\VerdierSourcePageContinue{0013}
\small
\begin{description}
\VerdierSemanticAnchor{front-publications-37}
\item[{[37]}] \emph{McKay 対応の幾何学的構成}（G. Gonzalez-Sprinberg
との共著）, Annales Scientifiques É.N.S. vol. 16, pp. 409--449,
Gauthier-Villars (1983).

\VerdierSemanticAnchor{front-publications-38}
\item[{[38]}] \emph{層の特殊化と穏やかなモノドロミー}, Analyse et
topologie sur les espaces singuliers (II--III), Astérisque 101--102,
pp. 332--364, S.M.F. (1983).

\VerdierSemanticAnchor{front-publications-39}
\item[{[39]}] \emph{幾何学的 Fourier 変換}（J.-L. Brylinski、
B. Malgrange との共著）, C.R.A.S. vol. 297, pp. 55--58,
Gauthier-Villars (1983).

\VerdierSemanticAnchor{front-publications-40}
\item[{[40]}] \emph{閉部分空間上への perverse 層の延長}, Systèmes
Différentiels et Singularités, Astérisque 130, pp. 210--217, S.M.F. (1985).

\VerdierSemanticAnchor{front-publications-41}
\item[{[41]}] \emph{モノドロミー的 perverse 層の延長}, Systèmes
Différentiels et Singularités, Astérisque 130, pp. 218--236, S.M.F. (1985).

\VerdierSemanticAnchor{front-publications-42}
\item[{[42]}] \emph{有理二重点上の反射加群}（M. Artin との共著）,
Math. Ann. 270, pp. 79--82, Springer (1985).

\VerdierSemanticAnchor{front-publications-43}
\item[{[43]}] \emph{正標数における McKay 則}（G. Gonzalez-Sprinberg
との共著）, C.R.A.S. vol. 301, pp. 585--587, Gauthier-Villars (1985).

\VerdierSemanticAnchor{front-publications-44}
\item[{[44]}] \emph{$S^2$ から $S^4$ への調和写像}, Geometry today,
International Conference, Rome 1984, Progress in Mathematics 60,
pp. 267--282, Birkhäuser (1985).

\VerdierSemanticAnchor{front-publications-45}
\item[{[45]}] \emph{有理二重点上の反射加群の乗法構造}
（G. Gonzalez-Sprinberg との共著）, Géométrie algébrique et applications:
I. Géométrie et calcul algébrique, Travaux en cours 22, pp. 79--110,
Hermann (1987).

\VerdierSemanticAnchor{front-publications-46}
\item[{[46]}] \emph{ある悲しい計算}, 付録, Géométrie algébrique et
applications: II. Singularités et géométrie complexe, Travaux en cours 23,
pp. 12--14, Hermann (1987).

\VerdierSemanticAnchor{front-publications-47}
\item[{[47]}] \emph{幾何学的 Fourier 変換 II}（J.-L. Brylinski、
B. Malgrange との共著）, C.R.A.S. vol. 303, pp. 193--198,
Gauthier-Villars (1986).

\VerdierSemanticAnchor{front-publications-48}
\item[{[48]}] \emph{量子群}（V. G. Drinfel'd による）, Séminaire
Bourbaki 1986--87, no. 685, Astérisque 152--153, pp. 305--319,
S.M.F. (1987).

\VerdierSemanticAnchor{front-publications-49}
\item[{[49]}] \emph{$S^2$ から $S^4$ への調和写像 II}, Harmonic mappings,
twistors, and $\sigma$-models, Advanced series in Mathematical Physics 4,
pp. 124--147, World Scientific (1988).
\end{description}
\normalsize
\VerdierSourcePageContinue{0014}
\small
\begin{description}
\VerdierSemanticAnchor{front-publications-50}
\item[{[50]}] \emph{新しい楕円ソリトン}, ``Algebraic Analysis'', Papers
Dedicated to Professor Mikio Sato on the Occasion of his Sixtieth Birthday,
Vol. II, pp. 901--910, Academic Press (1988).

\VerdierSemanticAnchor{front-publications-51}
\item[{[51]}] \emph{楕円ソリトン}（A. Treibich との共著）,
The Grothendieck Festschrift, Vol. III, Progress in Mathematics 88,
pp. 437--473, Birkhäuser (1990).

\VerdierSemanticAnchor{front-publications-52}
\item[{[52]}] \emph{接線被覆と4個の三角数の和}（A. Treibich との共著）,
C.R.A.S. vol. 311, pp. 51--54, Gauthier-Villars (1990).

\VerdierSemanticAnchor{front-publications-53}
\item[{[53]}] \emph{Boussinesq 楕円ソリトン：巡回の場合}
（E. Previato との共著）, Proceedings of the Indo-French Conference on
Geometry, pp. 173--185, Hindustan Book Agency (1993).

\VerdierSemanticAnchor{front-publications-54}
\item[{[54]}] \emph{K P 楕円ソリトンの Kritchever 多様体}
（A. Treibich との共著）, Proceedings of the Indo-French Conference on
Geometry, pp. 187--232, Hindustan Book Agency (1993).
\end{description}
\normalsize
\VerdierSourcePageSilent{0015}%
\clearpage
\begingroup
  \thispagestyle{empty}
  \centering
  \VerdierSemanticAnchor{thesis-title}
  \VerdierSourcePageBlock{0016}
  \vspace*{0.24\textheight}
  {\Huge アーベル圏の\par}
  \vspace{1.6em}
  {\Huge 導来圏\par}
  \vspace{7em}
  {\Large\bfseries J.-L. VERDIER\par}
  \vfill
\par
\endgroup
\VerdierSourcePageSilent{0017}%
